\chapter{Fundamentação Teórica}
\label{cap:fundamentacao}

% TODO(P1): Redigir a fundamentação teórica formal cobrindo os seguintes itens do Cap. 15 do Cormen:
% 1. Definição formal de Programação Dinâmica como método de projeto de algoritmos para problemas de otimização.
% 2. Propriedade 1: Subestrutura Ótima.
%    - Apresentar a demonstração formal pela técnica de "cortar e colar" (cut-and-paste) por contradição,
%      aplicada ao problema do Corte de Hastes (Cormen 15.1).
% 3. Propriedade 2: Superposição de Subproblemas (Overlapping Subproblems).
%    - Contrastar com divisão e conquista (onde os subproblemas são independentes).
%    - Grafo de subproblemas direcionado (DAG) e equivalência assintótica do número de vértices/arestas.
% 4. Os Quatro Passos do Projeto segundo Cormen:
%    - Passo 1: Caracterizar a estrutura de uma solução ótima.
%    - Passo 2: Definir recursivamente o valor de uma solução ótima.
%    - Passo 3: Calcular o valor de uma solução ótima (Bottom-Up vs. Top-Down com Memoização).
%    - Passo 4: Construir a solução ótima a partir das informações calculadas (reconstrução/traceback).
% 5. Comparativo Teórico: Programação Dinâmica vs. Algoritmos Gulosos vs. Divisão e Conquista.
% 6. Explosão Combinatória e Números de Catalan:
%    - Apresentar a fórmula P(n) = C(n-1) = (1/n) * \binom{2n-2}{n-1} para a cadeia de matrizes.
%    - Provar que P(n) = \Omega(4^n / n^{3/2}), justificando a necessidade incontornável da PD.

% [Texto a ser redigido pelo integrante P1 - ver docs/tarefas/P1_relatorio_secoes_1_2_6.md]
